% Apéndice: resolución de los ejercicios de afinidades
\section{Resolución de los ejercicios de afinidades}\label{cap:17}\label{ap:ejercicios-afinidades}

\subsection{Ejercicio 1: tres reflexiones puntuales}

\textbf{Enunciado.} Dados tres puntos no alineados \(a,b,c\), hallar el punto fijo de
\(\sigma_a\circ\sigma_b\circ\sigma_c\) y verificar que forma un paralelogramo con ellos.

\textbf{Solución algebraica.} La reflexión puntual de centro \(p\) es
\[
  \sigma_p(x)=2p-x.
\]
Entonces:
\[
\begin{aligned}
\sigma_a\sigma_b\sigma_c(x)
&=\sigma_a\!\left(\sigma_b(2c-x)\right)\\
&=\sigma_a(2b-2c+x)\\
&=2a-2b+2c-x\\
&=2(a-b+c)-x.
\end{aligned}
\]
La composición es una reflexión puntual de centro
\[
  \boxed{d=a-b+c}.
\]
Su punto fijo es \(d\), porque \(\sigma_d(d)=d\). Además:
\[
  b+d=a+c.
\]
Las diagonales \(ac\) y \(bd\) tienen el mismo punto medio, de modo que \(a,b,c,d\) forman
un paralelogramo.

\textbf{Interpretación geométrica.} Tomando \(b\) como origen provisional:
\[
  \overrightarrow{bd}
  =\overrightarrow{ba}+\overrightarrow{bc}.
\]
El punto \(d\) se obtiene mediante la regla del paralelogramo para sumar los vectores que van
desde \(b\) hacia \(a\) y hacia \(c\).

\subsection{Ejercicio 2: dos reflexiones puntuales}

\textbf{Enunciado.} Dados \(a\neq b\), demostrar que
\(\sigma_a\circ\sigma_b\) no tiene puntos fijos.

\textbf{Solución.}
\[
\sigma_a\sigma_b(x)=x+2(a-b).
\]
Si existiera un punto fijo:
\[
x+2(a-b)=x
\quad\Longrightarrow\quad
a=b.
\]
Esto contradice \(a\neq b\). Por lo tanto:
\[
  \boxed{\operatorname{Fix}(\sigma_a\sigma_b)=\varnothing
  \quad\text{si }a\neq b.}
\]

\textbf{Interpretación geométrica.} La composición es una traslación no nula. Todos los puntos
se desplazan por el mismo segmento orientado, por lo que ninguno puede permanecer en su lugar.
Si \(a=b\), las dos reflexiones se anulan y la composición es la identidad.

\subsection{Ejercicio 3: cuatro reflexiones en rectas}

\begin{atencion}[Ambigüedad del enunciado]
El texto comienza diciendo ``dadas cuatro rectas \(a,b,c,d\)'', pero termina llamando a
\(\varphi_i\) ``reflexión alrededor del punto \(i\)''. Por el contexto y la notación
\(\varphi\), se interpreta que son reflexiones respecto de las cuatro rectas.
\end{atencion}

Cada reflexión en una recta es una afinidad
\[
  \varphi_i(x)=H_ix+q_i,
\]
donde \(H_i\) es una matriz ortogonal de determinante \(-1\). La composición
\[
  F=\varphi_a\varphi_b\varphi_c\varphi_d
\]
tiene la forma
\[
  F(x)=Ax+t,
\]
con
\[
  A=H_aH_bH_cH_d,
\]
\[
  t=H_aH_bH_cq_d+H_aH_bq_c+H_aq_b+q_a.
\]
Como intervienen cuatro reflexiones,
\[
  \det A=(-1)^4=1.
\]
La transformación resultante preserva orientación y es una rotación o una traslación.

\textbf{Método analítico.} Los puntos fijos satisfacen
\[
  (I-A)x=t.
\]
\begin{itemize}
  \item Si \(I-A\) es invertible, existe un único punto fijo:
  \[
    \boxed{x_\ast=(I-A)^{-1}t}.
  \]
  La composición es una rotación con centro \(x_\ast\).
  \item Si \(A=I\) y \(t\neq0\), la composición es una traslación y no tiene punto fijo.
  \item Si \(A=I\) y \(t=0\), es la identidad y todos los puntos son fijos.
\end{itemize}

\textbf{Método gráfico.} Elegir dos puntos \(p,q\), construir sucesivamente sus cuatro
reflejos hasta obtener \(F(p)\) y \(F(q)\), y trazar las mediatrices de
\(\overline{pF(p)}\) y \(\overline{qF(q)}\). Si se cortan, la intersección es el centro de
rotación. Si son paralelas, el resultado es una traslación y no existe centro finito.

\subsection{Ejercicio 4: autovalor de una rotación en \(\R^3\)}

Sea \(R\in SO(3)\). Al ser una matriz real de dimensión impar, su polinomio característico
tiene al menos una raíz real. Como \(R\) es ortogonal, todo autovalor tiene módulo \(1\), por
lo que el autovalor real debe ser \(1\) o \(-1\).

Los otros dos autovalores, si no son reales, forman un par conjugado
\(e^{i\theta},e^{-i\theta}\), cuyo producto es \(1\). Como
\[
  \det R=1,
\]
el autovalor real restante debe ser
\[
  \boxed{\lambda=1}.
\]
Su autovector es la dirección del eje de rotación.

\begin{atencion}[Sobre la palabra ``único'']
Para una rotación no trivial de ángulo distinto de \(0\), \(1\) es el único autovalor real
asociado al eje, salvo casos degenerados. Para \(R=I\), todos los vectores son propios con
autovalor \(1\); por eso el enunciado necesita excluir la identidad si pretende unicidad del
eje.
\end{atencion}

\subsection{Ejercicio 5: matriz ortogonal que invierte orientación}

El enunciado afirma que una matriz ortogonal de \(\R^3\) con determinante \(-1\) tiene tres
autovalores reales iguales a \(1\) o \(-1\). Esa afirmación no es cierta en general.

Por ser real y de dimensión impar, la matriz posee al menos un autovalor real. Por ser
ortogonal, sus autovalores tienen módulo \(1\), y por tener determinante \(-1\), siempre
existe al menos un autovalor
\[
  \boxed{\lambda=-1}.
\]
Los otros dos pueden ser
\[
  e^{i\theta},\qquad e^{-i\theta},
\]
que no son reales cuando \(\theta\notin\{0,\pi\}\). Geométricamente, esto describe una
\textbf{rotorreflexión}: una rotación en un plano combinada con una inversión en la dirección
perpendicular.

\textbf{Contraejemplo.}
\[
Q=
\begin{pmatrix}
\cos\theta&-\sin\theta&0\\
\sin\theta& \cos\theta&0\\
0&0&-1
\end{pmatrix}
\]
es ortogonal y \(\det Q=-1\), pero sus autovalores son
\[
  e^{i\theta},\quad e^{-i\theta},\quad -1.
\]
Solo para casos especiales, como una reflexión en un plano, los tres son reales:
\[
  1,\ 1,\ -1.
\]

\subsection{Ejercicio 6: recuperar una afinidad desde tres puntos}

Sean \(a,b,c\) tres puntos no alineados y sean \(a',b',c'\) sus imágenes. Se busca
\[
  \gamma(x)=Ax+t.
\]
Restando la ecuación de \(a\) a las de \(b\) y \(c\):
\[
  b'-a'=A(b-a),
  \qquad
  c'-a'=A(c-a).
\]
Construimos las matrices cuyas columnas son los lados de ambos triángulos:
\[
  P=
  \begin{pmatrix}\vert&\vert\\b-a&c-a\\\vert&\vert\end{pmatrix},
  \qquad
  P'=
  \begin{pmatrix}\vert&\vert\\b'-a'&c'-a'\\\vert&\vert\end{pmatrix}.
\]
Como \(a,b,c\) no están alineados, \(P\) es invertible. Entonces:
\[
  \boxed{A=P'P^{-1}},
\]
y usando la correspondencia \(a\mapsto a'\):
\[
  \boxed{t=a'-Aa}.
\]
La transformación queda completamente determinada:
\[
  \boxed{\gamma(x)=P'P^{-1}x+a'-P'P^{-1}a}.
\]

\textbf{Comprobación de ausencia de puntos fijos.} Se intenta resolver
\[
  (I-A)x=t.
\]
Que la afinidad no tenga puntos fijos significa que este sistema es incompatible. En
coordenadas homogéneas:
\[
\begin{pmatrix}x'\\y'\\1\end{pmatrix}
=
\begin{pmatrix}
A_{11}&A_{12}&t_x\\
A_{21}&A_{22}&t_y\\
0&0&1
\end{pmatrix}
\begin{pmatrix}x\\y\\1\end{pmatrix}.
\]

\begin{resultado}[Estado de la guía de afinidades]
Los seis ejercicios quedan resueltos de forma algebraica y geométrica. Los ejercicios 3, 4
y 5 incluyen además las salvedades necesarias para interpretar correctamente sus enunciados.
\end{resultado}

